\documentclass[11pt]{article}
\usepackage[a4paper,margin=22mm]{geometry}
\usepackage{fontspec}
\setmainfont{LinLibertine_R.otf}[BoldFont=LinLibertine_RB.otf,ItalicFont=LinLibertine_RI.otf,BoldItalicFont=LinLibertine_RBI.otf]
\usepackage{xcolor}
\usepackage{hyperref}
\hypersetup{colorlinks=true,urlcolor=blue}
\setlength{\parindent}{0pt}
\setlength{\parskip}{6pt}
\sloppy
\emergencystretch=2em
\begin{document}
\section*{Historical Britannica: extracted source}
Encyclopaedia Britannica - Third edition, Volume 1, A-ANG - EB.5\par
{\small Source row 1801. Parser: nvidia/nemotron-3.5-parse.\par}
{\small This is model-extracted text, not a corrected transcription. OCR errors, reading order and region alignment require human review.\par}
\hrule\medskip
{\small\itshape 434\par}

{\scriptsize\color{gray} Region region-0; provider order 0.\par}
\medskip

\subsection*{**A L G E B R A.**}

{\scriptsize\color{gray} Region region-1; provider order 1.\par}
\medskip

{\small\itshape Part III.\par}

{\scriptsize\color{gray} Region region-2; provider order 2.\par}
\medskip

Application
to Gen
metry,
-006t nearly, and x=2.094 nearly.
\par

{\scriptsize\color{gray} Region region-3; provider order 3.\par}
\medskip

4. This operation may be continued to any length,
as by fuppeling gz=-054+/s, and fo ons, and the
value of z=0.094747 nearly.
\par

{\scriptsize\color{gray} Region region-4; provider order 4.\par}
\medskip

By the full operation a never value of x may be
found thus i force f=1 nearly and -1+1o/+of/+
i
f =0, f=0+of/r/s, that is, f=10+.6+.094
true to the left figure, and x=2.094.
\par

{\scriptsize\color{gray} Region region-5; provider order 5.\par}
\medskip

In the fame manner may the root of a pure equal Application
tion be found, and this gives an cafy method of ap-
proximating to the roots of numbers which are not
perfect powers.
\par

{\scriptsize\color{gray} Region region-6; provider order 6.\par}
\medskip

This can is applicable to numeral equations of every
order; and, by alluming a general equation, geneth
rules may be deduced for approximating to the roots
of any propried equation. By a finalime method we
may approximate to the roots of literal equations,
which will be exprefiled by infinite ferries.
\par

{\scriptsize\color{gray} Region region-7; provider order 7.\par}
\medskip

{\small\itshape P AR T III.\par}

{\scriptsize\color{gray} Region region-8; provider order 8.\par}
\medskip

{\small\itshape Of the Application of AlGEBRA TO GEOMETRY.\par}

{\scriptsize\color{gray} Region region-9; provider order 9.\par}
\medskip

{\small\itshape CHAP. I.\par}

{\scriptsize\color{gray} Region region-10; provider order 10.\par}
\medskip

{\small\itshape \_General Principles.\_\par}

{\scriptsize\color{gray} Region region-11; provider order 11.\par}
\medskip

C EOMETRY treats both of the magnitude and
polition of extention, and their connections.
\par

{\scriptsize\color{gray} Region region-12; provider order 12.\par}
\medskip

Algebra: tests only of magnitudely, therefore, of
the relations which fulfill in geometrical figures, those
of magnitude only can be immediately exprefed by
algebra.
\par

{\scriptsize\color{gray} Region region-13; provider order 13.\par}
\medskip

The opposite polition of friginal lines may indeed
be expelled firmly by the two + and — But, in
order to expect the various other positions of geomne-
tricical figures by algebra from the principles of geometry, forme relations of magnitude must be found, which
depend upon their path, positions, and which can be scalhi-
bited by equations ; And, convertily, by the fame
principles may the politions and figures is inferred from
the equations denoting fish relations of their parts.
\par

{\scriptsize\color{gray} Region region-14; provider order 14.\par}
\medskip

Though this application of algebra appears to be
indeted, yet fun as the simplicity of the operations,
and the general nature of its theorems, that investiga
tions, especially in the higher parts of geometry, are
generally order and more expeditions by the alge-
bracel method, though lef elegant than by what is
purely commented. The comentations albo, and ana-
legies of the two fiences established by this application, have given rifle to many curious tpecutions.
\par

{\scriptsize\color{gray} Region region-15; provider order 15.\par}
\medskip

Geometry has been rendered far more extenfire and
ufeful, and algebra itself.has received considerable im-
provements.
\par

{\scriptsize\color{gray} Region region-16; provider order 16.\par}
\medskip

{\small\itshape \#\# I. Of the Algebra Expertise of Geometrical Mag-
minder.\par}

{\scriptsize\color{gray} Region region-17; provider order 17.\par}
\medskip

A line, whether known or unknown, is reprefented
by a fragile letter: a \_relcangle\_ is properly exprefiled by
the product of the two letters reprefening its tides:
and a \_relcangular probability\_] by the product of three
factors ( wo of which represent the fields of any) its
rectangular bafes, and the third the altitude.
\par

{\scriptsize\color{gray} Region region-18; provider order 18.\par}
\medskip

There are the molth people copifications of geometrical
magnitudes, and any other which has a known pro-
portion to thefe, may in like manner be exprefled at
gebranically. Convertively the geometrical magnitudes,
reprefected by fich algebraic quantities, may be
found, only the global dimensions above the third,
one having any corresponding geometrical dimensions,
multi be exprefed by proportionals (4).
\par

{\scriptsize\color{gray} Region region-19; provider order 19.\par}
\medskip

The opposite polition of friginal lines, it has been
remained, may be exprefiled by the first + and—
\par

{\scriptsize\color{gray} Region region-20; provider order 20.\par}
\medskip

Thus, let a point A be given in the line
\par

{\scriptsize\color{gray} Region region-21; provider order 21.\par}
\medskip

AP, any fragment AP taken to the right hand being
considered as politive, a fragment Ap to the left is pro-
perly
\par

{\scriptsize\color{gray} Region region-22; provider order 22.\par}
\medskip

(a) All algebrical dimentions above the third multi be exprefiled by inferior geometrical dimensions; and
though any algebraic quantities of two or three dimensions may be immediately exprefiled by furfaces and
folds respectively, yet it is generally necefary to expref them, and till superior dimensions, by lines.
\par

{\scriptsize\color{gray} Region region-23; provider order 23.\par}
\medskip

If, in any geometrical investigation by algebras, each line is expelled by a finite letter, and each surface or
fold by an algebrical quantity of two or three dimensions respectively, then whatever legitimate operations are
performed with regard to them, the terms in any equation derived will, when properly reduced, be all of the
fame dimension, and any fuch equation may be cafily exprefilled geometrically by means of proportionals, as in
the following example.
\par

{\scriptsize\color{gray} Region region-24; provider order 24.\par}
\medskip

Thus, if the algebraic equation q-f+f=z=-pXI, to to be exprefed geometrically, a, f, o, and d, being fup-
pointed to repetis flingful here, let a1 e:cf/p, a contained proposition, then so f'd', c/g and a' c'+f+
b! :a:at+g; then let a:c:i:k;l, and a:c:c:a:l; alfo, let c:dim m:p, and c'd':e:p, or c' c'
--d':e:ce-p. By combining the two former proportions (Chap. II. Part I.), c':a'+b':::i:at+g, and
combining the latter with this lift found, c'--d':a'+b':e-pXI:cxat+g; therefore c-pXI=cxat+g, and
ecc-p:i:i:at+g.
\par

{\scriptsize\color{gray} Region region-25; provider order 25.\par}
\medskip

{\small\itshape IF\par}

{\scriptsize\color{gray} Region region-26; provider order 26.\par}
\medskip
\newpage\section*{Source record}
bri-1801-cb1efd626dfe\par
Image SHA-256: \texttt{cb1efd626dfe512b99ee297ecf32f9c5}\par\texttt{4d220d34efffac9442dbf2ed561195ac}\par
Dataset revision: \texttt{309f765b1ecc83e0a7f7}\texttt{762e44b4a927a7555443}\par
National Library of Scotland, Encyclopaedia Britannica Illustrated (CC0).\par
\url{https://huggingface.co/datasets/NationalLibraryOfScotland/encyclopaedia_britannica_illustrated/viewer/default/train?row=1801}
\end{document}
